\chapter{カスタマイズ --- GitHubを自分好みに整える}

\section{この章で学ぶこと}

道具は使いやすいように調整してこそ、その真価を発揮します。大工が鉋（かんな）の刃を研ぎ、料理人が包丁の柄を自分の手に馴染ませるように、開発者もGitやGitHubを自分の作業スタイルに合わせてカスタマイズすべきです。

この章では、GitHubとGitをより快適に使うためのカスタマイズ方法を学びます。

\begin{itemize}
  \item GitHubのテーマ設定（ダークモード/ライトモード）
  \item キーボードショートカットの活用
  \item GitHub CLI（gh）のエイリアス設定
  \item Gitコマンドのエイリアスと \cmd{.gitconfig} の書き方
  \item 通知の最適化 --- 情報の洪水に飲まれないために
  \item GitHub Mobileアプリの活用
  \item Saved Replies（保存済みの返信）による効率化
\end{itemize}

\section{テーマ設定 --- ダークモード/ライトモード}

長時間コードを読む開発者にとって、画面の明るさは目の疲れに直結します。GitHubでは複数のテーマが用意されており、好みや環境に合わせて選択できます。

設定方法は以下の通りです。

\begin{enumerate}
  \item 右上のプロフィールアイコンをクリック
  \item \textbf{Settings} を選択
  \item 左メニューの \textbf{Appearance} をクリック
  \item 好みのテーマを選択
\end{enumerate}

\screenshot{テーマ設定画面（Appearance）}

選択できるテーマは以下の通りです。

\begin{table}[h]
\centering
\begin{tabular}{ll}
\hline
テーマ & 特徴 \\
\hline
Light default & 標準的な白背景。明るい環境での作業に適している \\
Light high contrast & 白背景だがコントラストが高い。視認性を重視 \\
Dark default & 暗い背景。夜間の作業や暗い部屋での使用に適している \\
Dark high contrast & 暗い背景でコントラストが高い \\
Dark dimmed & やや明るめのダークテーマ。目に優しい \\
System に同期 & OSのダークモード設定に自動追従 \\
\hline
\end{tabular}
\end{table}

「System に同期」を選ぶと、macOSやWindowsのシステム設定でダークモード/ライトモードを切り替えたとき、GitHubのテーマも自動的に切り替わります。日中はライトモード、夜間はダークモードを使い分けたい場合に便利です。

\section{キーボードショートカット}

GitHubのWeb画面には多数のキーボードショートカットが用意されています。マウス操作を減らすことで、ページ間の移動やファイルの検索が格段に速くなります。まず覚えるべきは \textbf{\cmd{?} キー} です。どのページにいても \cmd{?} を押すと、そのページで使えるショートカットの一覧が表示されます。

\screenshot{?キーで表示されるショートカット一覧}

\subsection{グローバルショートカット（どのページでも使える）}

\begin{table}[h]
\centering
\begin{tabular}{ll}
\hline
ショートカット & 機能 \\
\hline
\cmd{?} & ショートカット一覧を表示 \\
\cmd{s} または \cmd{/} & 検索バーにフォーカスを移動 \\
\cmd{g} → \cmd{n} & 通知ページに移動 \\
\cmd{g} → \cmd{d} & ダッシュボード（ホーム画面）に移動 \\
\hline
\end{tabular}
\end{table}

\cmd{g} → \cmd{n} のように2つのキーを続けて押す操作は「コードシーケンス」と呼ばれます。最初の \cmd{g} を押した後、一定時間内に2つ目のキーを押す必要があります。\cmd{g} は "go to"（移動する）の略と考えると覚えやすいでしょう。

\subsection{リポジトリ内ショートカット}

リポジトリのページでは、各タブへの素早い移動が可能です。

\begin{table}[h]
\centering
\begin{tabular}{ll}
\hline
ショートカット & 機能 \\
\hline
\cmd{t} & ファイルファインダー（ファジー検索）を開く \\
\cmd{w} & ブランチ切り替えメニューを開く \\
\cmd{g} → \cmd{c} & Codeタブに移動 \\
\cmd{g} → \cmd{i} & Issuesタブに移動 \\
\cmd{g} → \cmd{p} & Pull requestsタブに移動 \\
\cmd{g} → \cmd{a} & Actionsタブに移動 \\
\cmd{g} → \cmd{b} & Projectsタブに移動 \\
\cmd{g} → \cmd{w} & Wikiタブに移動 \\
\hline
\end{tabular}
\end{table}

\subsection{ファイル表示時のショートカット}

\begin{table}[h]
\centering
\begin{tabular}{ll}
\hline
ショートカット & 機能 \\
\hline
\cmd{y} & パーマリンク（コミットハッシュ指定URL）に変換 \\
\cmd{l} & 行番号を入力してジャンプ \\
\cmd{b} & blameビュー（各行の変更者表示）に切り替え \\
\cmd{e} & ファイルを編集モードで開く \\
\hline
\end{tabular}
\end{table}

\subsection{Pull RequestとIssueのショートカット}

\begin{table}[h]
\centering
\begin{tabular}{ll}
\hline
ショートカット & 機能 \\
\hline
\cmd{r} & 選択テキストを引用して返信 \\
\cmd{Cmd/Ctrl + Enter} & コメントを送信 \\
\cmd{Cmd/Ctrl + Shift + p} & Markdownプレビューの切り替え \\
\hline
\end{tabular}
\end{table}

特に \cmd{Cmd/Ctrl + Shift + p} は、コメントを書いている途中でMarkdownがどのように表示されるかを確認するのに便利です。

\section{GitHub CLI (gh) のカスタマイズ --- エイリアスの設定と活用例}

GitHub CLI（\cmd{gh}）は第4章で紹介した通り、ターミナルからGitHubの操作ができるツールです。よく使うコマンドにはエイリアス（短縮名）を設定しておくと、タイプ量が減って作業が効率化されます。

\subsection{エイリアスの設定}

\begin{lstlisting}[language=bash]
# 基本的なエイリアス
gh alias set prc 'pr create'           # PR作成
gh alias set prv 'pr view --web'       # PRをブラウザで開く
gh alias set prl 'pr list'             # PR一覧
gh alias set il 'issue list'            # Issue一覧
gh alias set ic 'issue create'          # Issue作成

# 自分に関連するものだけ表示するエイリアス
gh alias set my-prs 'pr list --author @me'
gh alias set my-issues 'issue list --assignee @me'

# 現在のブランチのPRをチェックアウト
gh alias set prco 'pr checkout'

# 設定済みのエイリアス一覧を確認
gh alias list
\end{lstlisting}

設定したエイリアスは以下のように使います。

\begin{lstlisting}[language=bash]
# gh pr create の代わりに
gh prc --title "Add login feature"

# 自分が出したPRの一覧
gh my-prs

# 自分にアサインされたIssue一覧
gh my-issues
\end{lstlisting}

\subsection{ghの設定}

エイリアス以外にも、\cmd{gh config} コマンドでデフォルトの挙動をカスタマイズできます。

\begin{lstlisting}[language=bash]
# デフォルトエディタをVS Codeに設定
gh config set editor "code --wait"

# デフォルトブラウザを指定
gh config set browser "firefox"

# 対話的プロンプトを有効化（確認ダイアログを表示する）
gh config set prompt enabled

# 現在の設定を一覧表示
gh config list
\end{lstlisting}

\section{Gitのカスタマイズ --- エイリアスと.gitconfig}

Git自体にもエイリアス機能があります。\cmd{git status} を毎日何十回も打つ開発者にとって、\cmd{git st} と短縮できるだけでも積み重なると大きな時間の節約になります。

\subsection{エイリアスの設定}

\begin{lstlisting}[language=bash]
# 基本コマンドの短縮
git config --global alias.st status
git config --global alias.co checkout
git config --global alias.br branch
git config --global alias.ci commit
git config --global alias.sw switch

# ログ表示（見やすいフォーマット）
git config --global alias.lg "log --oneline --graph --decorate --all"
git config --global alias.ll "log --oneline -20"

# 差分の確認
git config --global alias.df diff
git config --global alias.ds "diff --staged"

# 直前のコミットにファイルを追加（メッセージは変えない）
git config --global alias.amend "commit --amend --no-edit"

# マージ済みブランチの一括削除
git config --global alias.cleanup "!git branch --merged main | grep -v main | xargs -n 1 git branch -d"
\end{lstlisting}

設定したエイリアスは即座に使えます。

\begin{lstlisting}[language=bash]
git st          # git status と同じ
git co main     # git checkout main と同じ
git lg          # きれいなグラフ付きログ
git ds          # ステージ済みの差分を確認
git amend       # 直前のコミットを修正
\end{lstlisting}

\cmd{git lg} は特におすすめのエイリアスです。ブランチの分岐と合流がグラフで視覚的に表示されるため、リポジトリの状態を直感的に把握できます。

\subsection{.gitconfig の直接編集}

エイリアスを一つずつ \cmd{git config} コマンドで設定するのが面倒な場合は、設定ファイルを直接編集することもできます。Gitのグローバル設定は \cmd{\textasciitilde/.gitconfig}（ホームディレクトリ直下の \cmd{.gitconfig}）に保存されています。

\begin{lstlisting}[language={}]
[user]
    name = Your Name
    email = your@email.com

[core]
    editor = code --wait
    autocrlf = input

[init]
    defaultBranch = main

[push]
    default = current

[pull]
    rebase = true

[alias]
    st = status
    co = checkout
    br = branch
    ci = commit
    sw = switch
    lg = log --oneline --graph --decorate --all
    ll = log --oneline -20
    df = diff
    ds = diff --staged
    amend = commit --amend --no-edit

[color]
    ui = auto

[diff]
    tool = vscode

[difftool "vscode"]
    cmd = code --wait --diff $LOCAL $REMOTE

[merge]
    tool = vscode

[mergetool "vscode"]
    cmd = code --wait $MERGED
\end{lstlisting}

いくつかの設定項目を補足します。

\begin{itemize}
  \item \textbf{\cmd{push.default = current}}: \cmd{git push} の際にブランチ名を省略すると、現在のブランチと同名のリモートブランチにプッシュする設定です。ブランチ名を毎回指定する手間が省けます。
  \item \textbf{\cmd{pull.rebase = true}}: \cmd{git pull} の際にマージコミットを作らず、リベースで履歴を直線的に保つ設定です。履歴をきれいに保ちたい場合に便利です。
  \item \textbf{\cmd{init.defaultBranch = main}}: \cmd{git init} で新しいリポジトリを作るときのデフォルトブランチ名を \cmd{main} にする設定です。
  \item \textbf{\cmd{color.ui = auto}}: ターミナルが対応していれば、diffやstatusの出力に色を付ける設定です。
\end{itemize}

\section{通知のカスタマイズ}

GitHubの通知は、設定を適切に行わないと大量のメールやアラートに埋もれてしまいます。「本当に重要な通知だけ受け取る」ように最適化しましょう。

\subsection{通知設定の場所}

Settings → \textbf{Notifications} で、通知の受け取り方を細かく設定できます。

\begin{lstlisting}[language={}]
Participating（参加中の会話）:
  Web通知:  有効にする
  メール:   有効にする
  → 自分が関わるIssue/PRの更新は確実に受け取りたい

Watching（ウォッチ中のリポジトリ）:
  Web通知:  有効にする
  メール:   無効にする（ノイズが多くなるため）
  → WebのNotificationsページで確認すれば十分

GitHub Actions:
  失敗時のみ通知にする
  → 成功したワークフローの通知は不要
\end{lstlisting}

\screenshot{通知設定画面}

\subsection{リポジトリごとの通知設定}

リポジトリのページ上部にある「\textbf{Watch}」ボタンで、リポジトリ単位の通知レベルを設定できます。

\begin{table}[h]
\centering
\begin{tabular}{ll}
\hline
設定 & 通知対象 \\
\hline
\textbf{Participating and @mentions} & 自分が参加中の会話と、@メンションされた場合のみ \\
\textbf{All Activity} & リポジトリ内のすべての活動（Issue作成、PR作成、コメントなど） \\
\textbf{Ignore} & このリポジトリからの通知を完全に無視 \\
\textbf{Custom} & Issue、PR、Release、Discussion などから受け取りたいものを選択 \\
\hline
\end{tabular}
\end{table}

おすすめの設定方針は以下の通りです。

\begin{itemize}
  \item \textbf{自分のリポジトリ}: All Activity（すべてを把握したい）
  \item \textbf{チームのリポジトリ}: Participating and @mentions（関係するものだけ）
  \item \textbf{参考にしているOSSリポジトリ}: Custom → Releases only（新バージョンだけ知りたい）
  \item \textbf{もう関わらないリポジトリ}: Ignore
\end{itemize}

通知が溜まりすぎたときは、GitHubの通知ページ（https://github.com/notifications）で「Mark all as read」を使って一括既読にできます。未読の通知を溜め込んでも何も良いことはありません。定期的にトリアージ（仕分け）する習慣をつけましょう。

\section{GitHub Mobile}

GitHub Mobileは、iOS/Androidに対応した公式アプリです。外出先でもスマートフォンからGitHubの主要な操作を行えます。

GitHub Mobileでできる主な操作は以下の通りです。

\begin{itemize}
  \item \textbf{Issue / PRの確認とコメント}: 通勤中や移動中にIssueの内容を確認し、コメントを返すことができます
  \item \textbf{コードの閲覧}: ファイルの中身を読むことができます（ただし編集はWeb版の方が快適です）
  \item \textbf{通知の管理}: プッシュ通知でリアルタイムに更新を受け取り、既読/未読の管理ができます
  \item \textbf{マージの承認}: Pull Requestのレビューを行い、Approveやマージができます
  \item \textbf{Discussionsへの参加}: ディスカッションの閲覧と投稿
\end{itemize}

特に便利なのは通知機能です。重要なPull Requestへのコメントや、CIの失敗通知をリアルタイムで受け取れるため、すぐに対応が必要な問題を見逃しにくくなります。

アプリのインストールは、App Store（iOS）またはGoogle Play Store（Android）から「GitHub」を検索してください。ログインにはWebブラウザ経由の認証が使われ、2FAにも対応しています。

\section{Saved Replies}

IssueやPull Requestのレビューでは、同じような内容のコメントを繰り返し書くことがあります。例えば、「LGTM（Looks Good To Me）」や「テストを追加してください」といったメッセージです。こうした定型文を毎回手入力するのは非効率です。

\textbf{Saved Replies（保存済みの返信）} は、よく使うコメントをテンプレートとして保存しておき、ワンクリックで挿入できる機能です。

\subsection{設定方法}

\begin{enumerate}
  \item \textbf{Settings} → \textbf{Saved replies} を開く
  \item テンプレートのタイトルと本文を入力して保存
\end{enumerate}

\subsection{便利なテンプレート例}

\begin{lstlisting}[language={}]
タイトル: LGTM
本文:
LGTM! Thank you for the contribution.

タイトル: Needs tests
本文:
Could you add tests for the new functionality?
We want to make sure this doesn't regress in the future.

タイトル: Duplicate
本文:
This appears to be a duplicate of #___. Closing in favor of the original issue.
Please add any additional context to the original issue.

タイトル: Stale
本文:
This issue has been inactive for a while. If the problem still exists,
please provide updated information. Otherwise, this will be closed.
\end{lstlisting}

Saved Repliesを使うには、IssueやPRのコメント欄の右上にあるメニューアイコンから選択するか、\cmd{Ctrl + .}（macOSでは \cmd{Cmd + .}）のショートカットで挿入できます。

チーム全体で同じSaved Repliesを使うことはできません（ユーザーごとの設定です）が、チーム内で「こういうテンプレートを使おう」というガイドラインを共有しておくと、コミュニケーションの一貫性が保たれます。

\section{まとめ}

この章では、GitHubとGitを自分好みにカスタマイズする方法を学びました。

\begin{itemize}
  \item \textbf{テーマ設定}: ダークモード、ライトモード、システム同期などから選べる。Appearance設定で変更
  \item \textbf{キーボードショートカット}: \cmd{?} キーで一覧表示。\cmd{t}（ファイル検索）、\cmd{g} → \cmd{i}（Issues移動）、\cmd{y}（パーマリンク変換）などが特に便利
  \item \textbf{GitHub CLI エイリアス}: \cmd{gh alias set} で頻用コマンドを短縮。\cmd{gh prc} で PR作成、\cmd{gh my-prs} で自分のPR一覧など
  \item \textbf{Git エイリアス}: \cmd{git config --global alias.st status} のように設定。\cmd{\textasciitilde/.gitconfig} を直接編集すると一括設定が楽
  \item \textbf{通知のカスタマイズ}: Participating はWeb+メール、Watching はWebのみ、Actions は失敗時のみが推奨設定。リポジトリごとのWatch設定で粒度を調整
  \item \textbf{GitHub Mobile}: iOS/Android対応。通知のリアルタイム受信とIssue/PRへのコメントが外出先でも可能
  \item \textbf{Saved Replies}: 定型コメントをテンプレート化。レビューの効率とコミュニケーションの一貫性が向上
\end{itemize}

カスタマイズに「正解」はありません。自分の作業スタイルに合わせて少しずつ調整し、快適な開発環境を作り上げていきましょう。
